\documentclass[11pt]{article}

\usepackage[T1]{fontenc}
\usepackage[margin=1in]{geometry}
\usepackage{amsmath}
\usepackage{amssymb}
\usepackage{booktabs}
\usepackage{array}
\usepackage{graphicx}
\usepackage{xcolor}
\usepackage{caption}
\usepackage{microtype}
\usepackage[colorlinks=true,linkcolor=blue!45!black,citecolor=blue!45!black,urlcolor=blue!45!black]{hyperref}

\hypersetup{pdftitle={FishAI v2.0: The Three-Sided Ask},pdfauthor={Allen Hsieh}}

% Float packing: let tables share pages with more text instead of drifting.
\renewcommand{\topfraction}{0.9}
\renewcommand{\bottomfraction}{0.6}
\renewcommand{\textfraction}{0.07}
\renewcommand{\floatpagefraction}{0.8}
\setcounter{totalnumber}{4}
\setcounter{topnumber}{3}

\newcommand{\code}[1]{\texttt{#1}}
\newcommand{\usff}{\code{us54}}

\title{FishAI v2.0: The Three-Sided Ask\\ Refuting an Avoidance Rule and Shipping Its Inverse}
\author{Allen Hsieh\\[2pt] {\small FishAI project}}
\date{August 2026}

\begin{document}
\maketitle

\begin{abstract}
\noindent Every ask in Literature (Canadian Fish) is three things at once: a bet on a card, a
\emph{broadcast} --- the holding licence of \usff{} row 6 makes an ask publish that the asker holds
a card of the named set --- and a \emph{concession}, because row 10 hands the turn to whoever was
asked on a miss. The engine has only ever scored the first. The project's owner asked for
``off-limits'' reasoning: identify the opponents who would punish a conceded turn, and refuse to ask
them. Implemented as a penalty on the miss branch and measured on duplicate deals against a control
that reproduces the shipped decision function byte-for-byte, it \textbf{loses} $4.5$ to $11.7$
points of win rate, monotonically worse as the appetite rises ($-0.52$ to $-1.61$ sets per duplicate
pair), while an information-free perturbation of the same magnitude is free ($-0.12 \pm 0.44$ at
matched size). The cause is structural and is this paper's central idea: row 6 makes threat and
opportunity \emph{the same fact}. A seat is dangerous exactly when it holds a card of a set this
team is invested in, which is what makes it the seat most likely to hold the card this team wants
--- $\mathrm{corr}(\text{threat}, \text{best available hit probability}) = +0.249$ over $33{,}595$
observations, mean best $p$ of $0.471$ at high-threat seats against $0.326$ elsewhere. The sound
response is therefore not avoidance but \textbf{defusal}: ask \emph{into} the set the dangerous seat
has published a basis in, take the card its reach rests on, and keep the turn under row 9. Same
threat model, opposite prescription, and it wins: $+1.43$ and $+1.58 \pm 0.32$ sets per duplicate
pair on two held-out banks, positive at all nine roster styles, $+1.08$ to $+1.96$ against opponents
that also defuse, and --- the validity result that matters most --- $+1.55$ to $+1.75$ inside a
third party's engine against their bots, with one knob moved. Most of the paper is negative.
Signalling is flat and its one positive arm failed to replicate ($-0.013 \pm 0.033$ at $1{,}500$
pairs); stalling is inert when weak and $-1.14 \pm 0.19$ sets when driven hard; the turn-handoff
request is built and refuted, the shipped hand-size rule already picking the best teammate $84.8\%$
of the time with $73.1\%$ of those decisions carrying no signal at all; and a contained-pass aim fix
measured below the instrument's floor and was \emph{reverted rather than shipped}, because a
listed-but-inert knob contaminates a measured vector. Concealment is a counter-weapon worth $+0.97
\pm 0.27$ added to defusal against defusing opponents, and \emph{worse than plain when played
alone}. This project's own intransitive-triangle claim is refuted at $4{,}000$ pairs per edge on two
independent bank sets --- the $2\times2$ is a strict total order $DC > D > P > C$ --- and replaced
by an opponent-dependent argmax worth one bit and up to $+0.29 \pm 0.14$. Two engine defects are
reported: \code{missTarget: 'fewest'} has the wrong sign (hand size correlates $-0.147$ with the
cards a conceded seat takes, against the licence estimator's $+0.421$), and \code{style.signalling}
is unreachable under \usff{}. Two facts are stated rather than buried: the certain-hit gate exists
because an adversarial pre-commit review found the defusal credit could displace a certain hit ($9$
occurrences over $150$ games $\times$ $9$ styles against $0$ for the control), and the project's own
detection-floor characterisation reclassifies two of this paper's nulls as non-measurements rather
than zeros.
\end{abstract}

\section{Introduction}
\label{sec:intro}

An ask in Literature is a question, and the engine has always treated it as one: rank the legal asks
by the probability that the named card is where you think it is, filter the long shots, break the
ties, play the argmax \cite{v05paper}. That model is correct as far as it goes, and it goes about a
third of the way. Under the \usff{} dialect \cite{rules} the same ask does two other things, both
consequential, neither scored. It \textbf{broadcasts}: row 6 permits an ask only into a set the
asker holds a card of and row 7 forbids naming a card you hold, so every ask publishes that the
asker holds a card of this set and lacks this one --- the only reliable public fact about a hidden
hand in this game, and what the engine's whole inference layer is built on \cite{obspaper,repo}. And
it \textbf{concedes}: row 9 keeps the turn on a hit, row 10 hands it to the seat that was asked on a
miss, so the miss branch of every ask carries a cost that depends on \emph{which} seat receives the
turn --- a cost the engine has priced by hand size alone.

The work began with a request from the project's owner, phrased as \emph{off-limits} reasoning:
build a model of which opponents would punish a conceded turn, and refuse to ask them. The threat
model that request describes is sound, derivable from the public log, and --- the first surprise ---
a strict improvement on what the engine already uses for the same job (Section~\ref{sec:cost}). The
prescription is not sound. Penalising asks by the danger of the seat they would concede to costs up
to $11.7$ points of win rate, and an information-free perturbation of the same size costs nothing,
so the loss is caused by the threat \emph{information} rather than by the disturbance
(Section~\ref{sec:offlimits}).

The reason is the paper's organising idea, and it fits in a sentence: \textbf{row 6 makes threat and
opportunity the same fact}. A seat is dangerous precisely when it holds a card of a set this team is
heavily invested in, and that is precisely what makes it the seat most likely to be holding the card
this team wants (Section~\ref{sec:confound}), so avoiding it means abandoning your own best set.

So keep the model and invert the prescription: ask the dangerous seat --- \textbf{defuse}, taking
back the very card its reach rests on while row 9 keeps the turn. That mechanism ships, and
Section~\ref{sec:defuse} measures it at increasingly hostile scales, ending inside a third party's
engine against a third party's bots (Section~\ref{sec:crossplay}). The rest is negative and
deliberately so: Sections~\ref{sec:signal} to~\ref{sec:conceal} report three further requested
mechanisms refuted and a fourth that works only as a counter to the first,
Section~\ref{sec:triangle} retracts a structural claim this project published about its own tactic
space, and Section~\ref{sec:floor} reports what the instrument underneath all of it can resolve
\cite{floorpaper}.

\section{Related work and prior results in this series}
\label{sec:related}

The measurement discipline is the project's own and is unchanged here: duplicate deals (every seeded
deal played in both orientations, so the deal is never a confound), the pair as the unit of
analysis, paired standard errors, and byte-exact controls wherever a mechanism can be switched off
\cite{v05paper,botlab}. Four prior results are load-bearing: the v0.5 roster's nine styles and tuned
Balanced control \cite{v05paper,styles}; the v1.0 counter table whose dominant row makes
style-switching provably degenerate, which Section~\ref{sec:capability4} prices the per-decision
layer against \cite{v10paper,v15paper,bounded}; the contained-book paper's value of a turn, $E =
\mathrm{hits}/\max(1,\mathrm{misses})$, and its independent estimate of the cards a conceded turn is
worth, which Section~\ref{sec:instrument} agrees with to within a few hundredths
\cite{containedpaper,containment}; and the detection-floor characterisation that supplies
Section~\ref{sec:floor}'s resolution limits \cite{floorpaper}. The external anchor is the
\textbf{fishlabs} repository \cite{fishlabs}, an independent Canadian Fish project with its own
engine, its own bots, and --- usefully --- an independently invented implementation of the very rule
Section~\ref{sec:offlimits} refutes.

\section{Background: the rules the results rest on}
\label{sec:background}

Six players in two teams of three (seats $0,2,4$ against $1,3,5$); 54 cards in nine half-suits
(``sets'') of six; nine cards dealt to each seat; the game ends the moment a team has been awarded
five sets, so ties are arithmetically impossible \cite{rules,pagat}. The rows this paper cites by
number \cite{rules}:

\begin{center}
\small
\begin{tabular}{@{}c p{0.85\textwidth}@{}}
\toprule
row & rule \\
\midrule
6  & \textbf{Ask licence.} The asker must hold at least one card of the asked set. \\
7  & \textbf{Own card.} You may not ask for a card you hold. \\
8  & \textbf{Target has cards.} The target must hold at least one card. \\
9  & \textbf{Hit.} The card transfers; the asker \emph{keeps the turn} and asks again. \\
10 & \textbf{Miss.} The turn passes to the player who was asked. \\
11 & \textbf{Declare timing.} Any time, including out of turn, inside the declare window. \\
14 & \textbf{Any declare error at all.} The opposing team scores the set. \\
17 & \textbf{Information.} Counts are public, with a full log of asks, results and declares. \\
\bottomrule
\end{tabular}
\end{center}

\noindent Rows 6 and 7 make the broadcast; rows 9 and 10 make the concession. Row 11 matters twice
below: it is why the value of a turn is the value of \emph{asks} rather than declares, and why the
turn-handoff request loses a premise (Section~\ref{sec:handoff}).

\section{What a conceded turn actually costs}
\label{sec:cost}

\subsection{The instrument}
\label{sec:instrument}

\code{scripts/probe-danger.mjs} records a concession whenever an ask is about to miss. The outcome
variable is how many cards the conceded seat then takes from the conceding team before it misses in
turn --- the length of row 9's hit chain, in cards --- read from the hands: this is a measurement
harness, not a bot. The \emph{estimate} it is compared against is built by \code{buildKnowledge}
over the acting seat's \code{SeatView}, i.e.\ exactly the information a policy has
\cite{concession,repo}. Over 300 \usff{} games, all-Balanced seats, $12{,}376$ conceded turns: the
mean cards taken per conceded turn is $\mathbf{1.304}$, with $P(0) = 46.7\%$, $P(\ge 2) = 33.2\%$
and a maximum of $11$. That mean lands within a few hundredths of the contained-book paper's
independent derivation of the same quantity from $E$ \cite{containedpaper,containment}, and two
arguments from different premises agreeing is the reason to trust either.

\subsection{Two estimators, one of which is backwards}
\label{sec:estimators}

\begin{itemize}
\item \textbf{\code{handSize}} --- \code{view.counts[X]}. What the engine uses today, in two places:
  \code{missTarget} in \code{decide.ts}'s \code{pickAsk}, and \code{valueContainedPass}'s $E \cdot
  n_t / \bar{n}$ term \cite{repo,containedpaper}.
\item \textbf{\code{licence}} --- over the sets $X$ has publicly shown a basis in (row 6, read off
  the log), the cards of those sets that this seat can certainly locate on its \emph{own} team.
  $X$'s reachable prey.
\end{itemize}

Over the $12{,}376$ observed concessions, \code{handSize} correlates $\mathbf{-0.147}$ with the
cards the conceded seat then took, and \code{licence} correlates $\mathbf{+0.421}$, with a
$3.2\times$ lift from its bottom decile ($0.778$ cards) to its top ($2.469$) \cite{concession}. That
is the first finding of the work, and it is a defect report: the shipped proxy has the wrong sign.
\code{missTarget: 'fewest'} prefers to concede to the opponent with the smallest hand, on the stated
grounds that the surrendered turn is worth least there; measured, the smallest hand takes
\emph{more}. Hand sizes fall as information accumulates, so the sign could have been a phase
artifact. It is not: over log lengths 0--40, 40--80 and 80--120 ($n = 5{,}216 / 4{,}626 / 2{,}452$,
mean cards taken $1.19 / 1.37 / 1.40$), $\mathrm{corr}(\code{hand})$ runs $-0.191 / -0.149 / -0.100$
against $\mathrm{corr}(\code{licence})$ at $+0.474 / +0.416 / +0.353$ --- negative in every bucket,
and Section~\ref{sec:handoff}'s independent fork study reproduces the sign a third time at $-0.261$.
Least squares over the licence estimator gives the coefficients the shipped \code{threat.ts}
carries:
\begin{equation}
\label{eq:threat}
\text{cards} \;\approx\; 0.885 \;+\; 0.391 \cdot \text{prey}.
\end{equation}

\subsection{Where the licence must be read from}
\label{sec:licence}

Not from \code{Knowledge.constraints}. The knowledge engine records an ask as a deal-time
set-constraint and \emph{drops it the moment it is satisfied or exhausted} --- which is exactly when
the seat has been shown to hold a card of the set, i.e.\ when the threat is most real. Reading the
basis off the public log instead multiplies what is visible, though not by the same factor twice:
detection of a harvest threat (opponent basis plus at least two of the set on my team) roughly
doubles, from $6.0\%$ to $\mathbf{11.7\%}$, while detection of the owner's strict five-and-one
position rises about fivefold, from $7.1\%$ to $\mathbf{35.0\%}$ \cite{concession}. Those figures are also why every mechanism here is a \textbf{graded term and
never a veto}: even at its best the predicate is invisible about two thirds of the time, and a rule
that fires on a third of its cases and is silent on the rest cannot be a hard prohibition without
being arbitrary.

\section{Off-limits, as requested: refuted}
\label{sec:offlimits}

The request, implemented literally, is a penalty on the miss branch,
\begin{equation}
\label{eq:penalty}
\text{penalty} \;=\; w_{\mathrm{hit}} \cdot (1-p) \cdot a \cdot \frac{D(t)}{1+E},
\end{equation}
subtracted from the shipped ask score, where $p$ is the seat's own hit estimate, $D(t)$ the threat
estimate of Equation~\eqref{eq:threat} at target $t$, $E$ the value of a turn in cards, and $a$ the
appetite. Two hundred duplicate pairs per cell, against a prototype control that reproduces
\code{decide} byte-for-byte when the mechanism is off \cite{concession}.

\begin{table}[t]
\centering
\caption{The owner's off-limits rule swept over appetite: a danger penalty on the miss
branch, 200 duplicate pairs per cell \cite{concession}. The first column is the byte-exact control.}
\label{tab:offlimits}
\footnotesize
\begin{tabular}{@{}lrrrrrr@{}}
\toprule
appetite & $0$ (control) & $0.05$ & $0.10$ & $0.25$ & $0.50$ & $1.00$ \\
\midrule
win rate & $50.00\%$ & $45.0\%$ & $43.0\%$ & $45.5\%$ & $42.3\%$ & $\mathbf{38.3\%}$ \\
paired set-difference & $0.0000 \pm 0.0000$ & $-0.52$ & $-0.73$ & $-0.89$ & $-1.19$ & $\mathbf{-1.61}$ \\
\bottomrule
\end{tabular}
\end{table}

Table~\ref{tab:offlimits} is monotone in the wrong direction: the harder the rule is driven, the
worse it does, from $-0.52$ to $-1.61$ sets per duplicate pair and $4.5$ to $11.7$ points of win
rate below the control. A hard veto restricted to the narrow case --- likely-miss asks only, at
targets with at least $K$ certainly-located prey cards --- loses too, at $-1.02$ to $-1.38$ for $K =
2 \ldots 5$, reaching neutral only at $K = 6$, where it has stopped firing.

\subsection{The null control, which is what makes this a finding}
\label{sec:null}

A policy perturbation that costs win rate has two candidate causes: the information it uses, or the
mere fact of disturbance. The separating experiment is the same penalty at the same magnitude with
$D$ replaced by an information-free hash of the target. At appetites $0.10 / 0.25 / 0.50$ the null
measures $48.50\% / 50.25\% / 50.00\%$ win rate and $-0.12 \pm 0.44$, $+0.01 \pm 0.42$, $-0.23 \pm
0.44$ sets per pair \cite{concession}. \textbf{An information-free perturbation of the same size is
free}, and the danger-informed one costs $4.8$ to $7.7$ points more at matched magnitude ($48.5$
against $43.0$; $50.25$ against $45.5$; $50.0$ against $42.3$). The loss is caused by the threat
information, not by the noise.

\subsection{Why: threat and opportunity are the same fact}
\label{sec:confound}

Row 6 cuts both ways, and here is what that means numerically. Over $33{,}595$ ask decisions,
$\mathrm{corr}(\text{threat}(t), \text{best hit probability available at } t) = \mathbf{+0.249}$:
the $12{,}691$ high-threat targets (at least two reachable prey cards) carry a mean best available
$p$ of $\mathbf{0.471}$ against $\mathbf{0.326}$ at the $20{,}904$ low-threat ones
\cite{concession}. A seat is dangerous \emph{because} it holds a card of a set this team is heavily
invested in --- and that is the same fact that makes it the seat most likely to be holding the card
this team wants. Penalising threat is penalising your own best set.

There is also a soundness argument that reaches the same conclusion without any win rate, and it
generalises past this engine. A threat model built from public information is a strict
\textbf{under-approximation} of what a seat knows. A positive fact in it --- \emph{this seat has
published a basis in that set} --- is a certificate and is sound. An absence is not evidence: the
seat may hold a basis it has never had occasion to publish. So the model may legitimately be used to
act on \emph{proven} reach, and may \textbf{never} be used to certify a seat as safe. Avoidance
requires exactly the forbidden direction; defusal requires only the sound one.

\section{Defusal: the same model, the opposite prescription}
\label{sec:defuse}

If the dangerous seat is the seat holding the card you want, the answer is not to leave it alone ---
it is to take the card back. A hit removes the very card the licence rests on, and row 9 keeps the
turn while doing it. The mechanism is therefore a bonus on the \emph{hit} branch where the
off-limits rule was a penalty on the \emph{miss} branch:
\begin{equation}
\label{eq:defuse}
\text{bonus} \;=\;
\begin{cases}
a_{d} \cdot w_{\mathrm{hit}} \cdot p \cdot \dfrac{c \cdot \mathrm{prey}(B)}{1+E}
  & \text{if the target has a published basis in } B,\\[2ex]
0 & \text{otherwise,}
\end{cases}
\end{equation}
added to the ranked ask score, where $B$ is the asked card's set, $\mathrm{prey}(B)$ the number of
cards of $B$ certainly located on this team, $c$ the fitted cards-per-prey-card slope of
Equation~\eqref{eq:threat}, and $a_d$ the style's appetite \cite{repo}. Dividing by $1+E$ converts
the credit from cards into the ranker's own units, so the appetite means the same thing at every hit
rate.

Two design decisions were measured rather than assumed. \textbf{Per set, not per seat}: crediting
the target's whole published reach was measured too, and per-set is both better ($+1.50$ against
$+1.21$ on the same bank) and the only version whose causal story is true, since a hit on a card of
$B$ cannot protect a set it is not in. \textbf{A term, not a branch}: because the credit is added to
the score, two of \code{pickAsk}'s three guarantees survive for free --- \code{minHitP} still
filters the same pool, and both near-tie windows still break the same ties. The third did not.

\subsection{The certain-hit gate, and why it exists}
\label{sec:gate}

\textbf{The gate was not designed in. It was added because an adversarial pre-commit review found
that the credit could displace a certain hit, and then measured how often.} The arithmetic is thin
in the wrong direction: a certain hit's score margin over an uncertain ask is about $2$ points,
while Equation~\eqref{eq:defuse} is bounded only by $a_d \cdot w_{\mathrm{hit}} \cdot c \cdot
\mathrm{prey} / (1+E)$, which reaches about $137$ at $w_{\mathrm{hit}} = 70$. Ungated, the term
abandoned a certain hit \textbf{9 times over 150 games $\times$ 9 styles, against 0 for the control}
\cite{concession,repo}.

\code{pickAsk} therefore zeroes \emph{both} concession terms for any ask with $p < 1$ whenever a
certain hit is available, leaving the credit live \emph{among} certain hits, where it is free.
Section~\ref{sec:shipped}'s table was re-measured with the gate in place and did not move. Two
honest riders travel with it, both in Section~\ref{sec:open}: its firing rate is unpriced beyond the
sweep that found it, and its concealment half is not pinned by a test.

\subsection{Measured, on held-out banks and across the roster}
\label{sec:shipped}

The appetite was fixed at $1$ on a tuning bank --- the sweep is an inverted $U$ peaking at $1$--$2$
and back to noise by $8$ --- then confirmed on two disjoint held-out banks and across the roster.
Table~\ref{tab:defuse} is the \emph{shipped} implementation, not the prototype:
\code{STYLE\_ROSTER[s]} against the same style with \code{defuse: 0}, 400 duplicate pairs per cell
\cite{concession}.

\begin{table}[t]
\centering
\caption{Defusal, shipped implementation: each roster style against itself with
\code{defuse: 0}, 400 duplicate pairs per cell \cite{concession}. Intervals are $\pm 1.96$ paired
SE.}
\label{tab:defuse}
\footnotesize
\begin{tabular}{@{}lrr@{}}
\toprule
cell & win rate & paired set-difference \\
\midrule
balanced, holdout A & $60.88\%$ & $\mathbf{+1.43 \pm 0.32}$ \\
balanced, holdout B & $61.38\%$ & $\mathbf{+1.58 \pm 0.32}$ \\
blitz     & $61.88\%$ & $+1.35 \pm 0.32$ \\
punter    & $58.38\%$ & $+1.22 \pm 0.33$ \\
banker    & $57.63\%$ & $+1.12 \pm 0.31$ \\
turtle    & $61.25\%$ & $+1.35 \pm 0.36$ \\
hoarder   & $59.75\%$ & $+1.31 \pm 0.35$ \\
scout     & $61.63\%$ & $+1.66 \pm 0.32$ \\
ghost     & $55.38\%$ & $+0.83 \pm 0.32$ \\
archivist & $\mathbf{65.00\%}$ & $\mathbf{+1.94 \pm 0.33}$ \\
\bottomrule
\end{tabular}
\end{table}

Every interval excludes zero; the sign replicates on both banks and at all nine styles.

\textbf{Two sets of numbers appear in this literature and they are not the same measurement.} The
prototype --- the term alone, hooked into a mirror of \code{pickAsk} --- measured $+1.50\ [1.24,
1.75]$ and $+1.65\ [1.38, 1.91]$ on 600-pair held-out banks, and those are the figures
\code{defuse.ts} carries in its header as the result it was built from. Table~\ref{tab:defuse} is
the shipped implementation re-measured at 400 pairs and lands at $+1.43$ and $+1.58$ on the same two
banks. The gap is inside the intervals and is what a smaller bank plus a real integration should
look like; both are quoted rather than the flattering one. Three limits travel with the table, two
of which the next sections discharge. It is a \emph{mirror match}, the arena in which a mechanism is
most flattered (Section~\ref{sec:gauntlet}). The slope $c$ is \emph{fitted on this engine}, against
this roster, under \usff{} --- one fitted constant, and the only one in the mechanism
(Section~\ref{sec:crossplay}). And the prototype touched \emph{one hook}: nothing in the declare
policy, the containment policy or the adaptive layer was changed.

\subsection{The gauntlet: against opponents that also defuse}
\label{sec:gauntlet}

The fishlabs corpus supplies the warning that motivates this experiment: their structurally
analogous inverted coordinate measured $+2.71$ against a weak opponent, $+0.55$ against a middling
one and $-0.99$ against their strongest --- a gain that existed only against opponents unable to
punish it. Those three figures are recorded in this project's own concession document rather than in
the cross-play one \cite{concession}. So the gate is whether defusal still pays when the opponent
defuses too. Both arms face the \emph{same fully-armed opponent} and only the measured team's
appetite changes: arm~A is \code{balanced(defuse 1)} against $S(\code{defuse 1})$, arm~B
\code{balanced(defuse 0)} against the same opponent, 250 duplicate pairs per cell on a third bank,
reported as A minus B.

\begin{table}[t]
\centering
\caption{The gauntlet: both arms against the same fully-armed opponent, 250 duplicate pairs
per cell \cite{concession}. Arm~A against \code{balanced} is two identical policies on duplicate
deals and must sum to zero by construction.}
\label{tab:gauntlet}
\footnotesize
\begin{tabular}{@{}lrrr@{}}
\toprule
opponent style & arm A & arm B & \textbf{delta} \\
\midrule
balanced  & $0.000$  & $-1.636$ & $\mathbf{+1.64 \pm 0.40}$ \\
blitz     & $-0.028$ & $-1.744$ & $\mathbf{+1.72 \pm 0.45}$ \\
punter    & $-0.228$ & $-2.156$ & $\mathbf{+1.93 \pm 0.45}$ \\
banker    & $+0.436$ & $-1.112$ & $\mathbf{+1.55 \pm 0.55}$ \\
turtle    & $+3.304$ & $+1.724$ & $\mathbf{+1.58 \pm 0.56}$ \\
hoarder   & $+1.300$ & $-0.584$ & $\mathbf{+1.88 \pm 0.46}$ \\
scout     & $+1.876$ & $-0.088$ & $\mathbf{+1.96 \pm 0.51}$ \\
ghost     & $+0.580$ & $-0.504$ & $\mathbf{+1.08 \pm 0.52}$ \\
archivist & $+0.696$ & $-0.916$ & $\mathbf{+1.61 \pm 0.53}$ \\
\bottomrule
\end{tabular}
\end{table}

The gain survives at every style and is if anything larger than in the mirror match. Three features
deserve reading separately. \textbf{Arm~A against \code{balanced} is exactly $0.000$}: two identical
policies on duplicate deals must sum to zero by construction, and they do to the last digit --- a
free internal-validity check on the harness, and it passes. \textbf{Arm~B is negative almost
everywhere}: a team that does not defuse \emph{loses} to one that does, by roughly the margin it
would have gained by switching on, so the mechanism is contested rather than a coordination gain
both sides collect, and declining to play it is a position rather than a neutrality. And
\textbf{\code{ghost} is the narrowest cell} ($+1.08$), the style with the most suppressed
information use and therefore the least to defuse with --- an ordering that is mechanistic, and mild
evidence that the effect is the stated one. This is not unconditional --- every opponent here is
still a FishAI policy, and the strongest \emph{known} opponent to this engine need not be the
strongest one --- but the specific failure the corpus predicted did not occur.

\section{The cross-engine replication}
\label{sec:crossplay}

Section~\ref{sec:gauntlet} answers ``does it survive a stronger opponent'' without answering ``is
this a fact about Canadian Fish or a fact about this roster''. The second question needs a different
engine, and there is one: FishAI's \code{decide} was run as a guest inside the \textbf{fishlabs}
TypeScript engine, against \emph{its} bots, with \code{defuse} the only knob moved \cite{crossplay}.

\subsection{What could not be played, stated first}
\label{sec:sestina}

The owner's request was for simulations against \textbf{SESTINA v1.0}, that project's frontier
agent. \textbf{No game below is a game against SESTINA v1.0.} It is C++ and this machine has no C++
toolchain --- \code{g++}, \code{clang++}, \code{cl}, \code{gcc}, \code{make} and \code{cmake} are
all absent --- so their binary cannot be produced here. What was played instead is their
\emph{TypeScript} lab engine and the eight bots shipped in it, which run under Node with no
compiler. Their strongest TypeScript bot is FishBot v0.3; their lineage runs to v0.7, which is
SESTINA v1.0, so the opponent here is roughly four releases behind their frontier \cite{crossplay}.
One naming correction travels with that: the agent was briefly called ``SISTINA'', a first search
returned nothing, and it momentarily looked as though it did not exist. It does, under the other
spelling \cite{fishlabs}.

\subsection{The bridge, and the handicap it imposes}
\label{sec:bridge}

A bridge that quietly changes the game measures nothing, so the two rule sets were compared before
any number was trusted. Deck, seats and teams, ask legality, turn flow and the wrong-declare outcome
are all identical between the fishlabs TypeScript engine and \usff{}, which makes the adapter a pure
renaming and means FishAI's inference layer reads a log that means exactly what it means at home.
One row differs: in their engine \textbf{only the turn-holder may declare}, where \usff{} row 11
offers every seat the option in the declare window \cite{crossplay}.

That difference handicaps FishAI, and it was quantified rather than waved at. Measured in FishAI's
own engine over 300 \usff{} games, $54.6\%$ of its declares are made on its own turn --- the channel
the host provides --- and $\mathbf{45.4\%}$ off-turn, through a door their engine does not have.
Some of those sets would simply be declared later on the seat's own turn, so $45.4\%$ is an upper
bound on the loss rather than an estimate of it; but the direction is unambiguous, and every FishAI
win rate in the host engine is a \textbf{floor}. The adapter fell back to a default in 6 decisions
out of roughly nine thousand ($0.07\%$), always on the declare-window move that does not exist
there, and results are bit-identical with and without it.

\subsection{Standing, and then the result that matters}
\label{sec:standing}

\begin{table}[t]
\centering
\caption{Defusal inside the fishlabs engine against the fishlabs bots, 300 duplicate pairs
per opponent, \code{defuse} the only knob moved \cite{crossplay}.}
\label{tab:xplay}
\footnotesize
\begin{tabular}{@{}lrrr@{}}
\toprule
opponent & \code{defuse: 1} net sets & \code{defuse: 0} net sets & \textbf{paired delta} \\
\midrule
\code{fishbot} (v0.3)  & $-1.333$ & $-2.880$ & $\mathbf{+1.547 \pm 0.517}$ \\
\code{fishbot\_v02}    & $-0.933$ & $-2.627$ & $\mathbf{+1.693 \pm 0.555}$ \\
\code{lockout}         & $-0.573$ & $-2.327$ & $\mathbf{+1.753 \pm 0.533}$ \\
\code{detective}       & $-0.593$ & $-2.300$ & $\mathbf{+1.707 \pm 0.560}$ \\
\bottomrule
\end{tabular}
\end{table}

First, standing. At 150 games per cell in both orientations, FishAI at style \code{balanced} wins
$42.67\%$ against \code{fishbot} (v0.3), $45.33\%$ against \code{detective}, $47.00\%$ against
\code{lockout} and $48.67\%$ against \code{fishbot\_v02}, and $80.67$ to $99.67\%$ against the other
four bots in their lab \cite{crossplay}: competitive with their top four, one to seven points under
even against them, while playing without the channel that carries $45\%$ of its declares. Those four
cells publish no intervals and Section~\ref{sec:floor} reclassifies them accordingly.

Table~\ref{tab:xplay} is the strongest evidence in this project that defusal is a property of the
game rather than of FishAI's roster. Four independent opponents, four intervals excluding zero ---
and the effect \emph{size} transfers: $+1.55$ to $+1.75$ here against $+1.50$ and $+1.65$ on
FishAI's own held-out banks and $+1.08$ to $+1.96$ against FishAI opponents that also defuse. The
mechanism was derived, fitted and tuned entirely against FishAI's roster and had never seen any of
these bots. Note what the table does \emph{not} say: the \code{defuse: 1} column is still negative
against all four. Defusal recovers about $1.6$ sets of a deficit it does not close; FishAI loses to
their top bots either way, by less with the mechanism on.

Since SESTINA could not be run, their reported evidence is cited as theirs
\cite{fishlabs,crossplay}: SESTINA v1.0 beats their pre-registered target by $+3.33$ pp $[+2.88,
+3.78]$ over $48{,}000$ games and the deployed v0.6 policy by $+4.63$ pp $[+4.19, +5.06]$ over the
same $48{,}000$. \emph{Frontier margins in this game are small} --- a seventh development cycle with
a determinized search over a particle belief buys about three percentage points --- which puts
$+1.6$ sets per duplicate pair in perspective as a large effect by this game's standards.

Their README pairs those margins with three qualifications, and quoting the margins without them
would overstate the frontier this paper measures itself against. They report that SESTINA costs
\textbf{at least $4.52\times$} the non-searching blueprint. They publish a negative about their own
flagship: it does \emph{not} measurably beat a composite configuration assembled earlier in the same
development cycle. And they decline the strongest claim outright, separating ``strongest
configuration this project has produced'', which they record as established, from global standing,
which they record as \emph{``not established, not claimed''} --- so there is no published basis,
theirs or this paper's, for treating SESTINA as the strongest Canadian Fish agent in existence
\cite{fishlabs,crossplay}.

\subsection{\code{lockout}: the owner's rule, built independently, and doing well}
\label{sec:lockout}

The tension in this paper's central negative result is in their repository, and it is recorded here
rather than buried. Their strategy table describes \code{lockout} as a turn-starvation specialist
that ranks asks by posterior and then \emph{charges a penalty for missing into a dangerous opponent}
--- the owner's off-limits rule, arrived at independently by someone else. And in their own
round-robin (every fishlabs bot against every other, 100 games $\times$ 2 orientations per pairing)
it ranks second of nine at $77.06\%$ mean win rate against the field, behind \code{fishbot} v0.3's
$82.69\%$ and ahead of \code{fishbot\_v02} ($75.50\%$) and \code{detective} ($75.31\%$)
\cite{crossplay,fishlabs}.

\textbf{This looks like a contradiction of Section~\ref{sec:offlimits} and it is not}, for a reason
that separates a strategy from a term. \code{lockout} is a whole strategy row --- its own
declaration threshold, its own risk parameter, its own targeting --- that happens to \emph{contain}
a danger penalty. Its $77\%$ says a strategy built that way is strong; it does not isolate the
penalty. Section~\ref{sec:offlimits} is an ablation: one policy, one term added, everything else
byte-identical, duplicate deals, plus an information-free null at matched magnitude. Only the
ablation answers ``does the danger penalty help''.

\textbf{A supporting argument offered here has been withdrawn.} An earlier draft read the $34.1\%$
ask accuracy recorded in the \code{lockout} cell as \code{lockout}'s own, and took it as this
paper's mechanism visible from outside --- a policy that avoids its best asks showing a low hit
rate. That was a misread of a column that is FishAI's. Measured properly, both sides ask at almost
exactly the same rate there ($34.1\%$ against $34.8\%$; and $49.8/50.6$, $45.4/46.0$, $43.5/44.4$ in
the other three cells), so the low rate is a property of the \emph{matchup}, depressing both sides
alike \cite{crossplay}. The inference is retracted rather than restated; what survives is the
narrower observation that \code{lockout} produces the largest \code{defuse} delta of the four
($+1.753$), consistent with an opponent that publishes bases freely but establishing nothing about
why. So the tension is left unresolved, which is the honest state, and only an ablation of
\code{lockout}'s own coefficient inside their engine would settle it --- which their lab does not
expose as a knob.

\subsection{Common failure points, theirs and this paper's}
\label{sec:failures}

The owner's request named common failure points explicitly, and the host lab answers it most sharply
through its own weakest serious bot. \code{kv\_search} collapses on \emph{declaration calibration},
not on play. It takes $30.25\%$ against the field and $0.0\%$ against each of the top three, and
over 40 games against \code{fishbot} it takes $1.5$ sets to $7.5$ --- but its ask accuracy in those
games is $49.5\%$ against $62.1\%$, a gap far too small to carry that result. Its declaration
accuracy does carry it: $\mathbf{73.0\%}$ against \code{fishbot}'s $\mathbf{95.3\%}$, from the
lowest declaration threshold of any serious bot in their table ($0.75$) \cite{crossplay}. It is not
being out-asked into oblivion; it is \emph{gifting sets}. Under a rule set where any declare error
awards the set to the opposing team --- \usff{} row 14, and their engine agrees
(Section~\ref{sec:bridge}) --- a $27\%$ misdeclare rate is a losing position no amount of asking
recovers. The lesson generalises beyond their bot: \textbf{under gift rules, declare calibration
dominates ask quality.}

It also lands on this paper, which is the most self-critical thing the cross-play produced. FishAI's
own declare accuracy against their bots is $90$--$94\%$, against \code{fishbot}'s $97\%$, and that
gap is a plausible share of the one-to-seven points FishAI sits below their top four in
Section~\ref{sec:standing} \cite{crossplay}. Its other visible weakness in this host is the declare
channel of Section~\ref{sec:bridge}: without the off-turn window it must hold a located set until
its own turn comes round, and a set held is a set an opponent may take a card out of first.

\section{Signalling: refuted, and its observational support was reverse-causal}
\label{sec:signal}

An ask publishes two facts to \emph{teammates} as well as to opponents. The hypothesis is natural
and its observational support is strong: sets that are declared correctly have been asked into far
more often than sets declared wrongly. \code{scripts/probe-declare.mjs}, 200 games per style, every
declare classified by whether the ask channel was still open --- i.e.\ whether row 8 still permitted
an ask at all --- gives, for Balanced and Punter respectively: $1{,}468$ and $1{,}469$ declares with
the channel open, $\mathbf{100.0\%}$ and $\mathbf{99.7\%}$ of them correct; $90$ and $84$ with it
closed, $\mathbf{78.9\%}$ and $\mathbf{75.0\%}$ correct; and a mean of $11.70$ / $11.80$ prior asks
into the set before a \emph{correct} declare against $\mathbf{7.05}$ / $\mathbf{7.28}$ before a
wrong one \cite{concession}.

\textbf{The intervention did not reproduce the association.} A bonus for asking into sets the team
cannot yet prove, weighted by proximity to channel-close, measured $+0.013$, $+0.023$, $+0.067$,
$-0.090$, $-0.203$ at rising appetite over 300 pairs --- and the one positive arm \textbf{failed to
replicate}: $1{,}500$ duplicate pairs on a disjoint bank gave $49.97\%$, $-0.013 \pm 0.033$. The
likely reading is that the correlation is reverse-causal: a set asked into a lot is a set several
seats hold a basis in, and \emph{that} is what makes it provable, so the ask count is a symptom of
provability rather than a cause of it. A second version is also null --- spending only the
\emph{free} choice, a tie-break among near-equal asks preferring the one that most reduces the
team's allocation ambiguity at zero cost in hit probability, measured $+0.008 \pm 0.074$ and $-0.010
\pm 0.069$ on two banks.

\subsection{The structural objection, which is bigger than signalling}
\label{sec:structural}

Two things this refutation does not rule out, and the first governs Section~\ref{sec:stall} as well.

\textbf{A structural objection to the whole family.} A linear ask score whose positive mass is
hit-probability-driven cannot express a \emph{deliberate miss}: the hit term is zero by construction
and every remaining term is a penalty, so such an ask can never win an argmax at any weight.
FishAI's score has exactly that shape --- set $p$ to $0$ and $70$ of the mass vanishes. So
``signalling is inert when weak'' may be a fact about the scoring function rather than about
signalling; the measured \emph{harm} at high weight is real, but the inert regime is precisely what
a structurally unselectable tactic looks like. Defusal escapes the trap only because a defusal ask
\emph{is a hit}, and concealment (Section~\ref{sec:conceal}) only because it moves the argmax
between two hits.

\textbf{The receiver side is untested.} Everything above changes what the \emph{sender} asks. A
teammate that reads the log's row-6 and row-7 facts more aggressively when forming a declare plan is
a different experiment, and has not been measured.

\section{Stalling: refuted}
\label{sec:stall}

The owner's third request was an endgame stall: decline to take the opponents' last card, to keep
the ask channel open. Under \usff{} there is no endgame phase --- the rule set specifies
\code{wholeTeamOut: 'declareWindow'} and the reducer simply moves the turn to the next seat with
cards. What emptying the opponents does is end the \textbf{ask channel}, since row 8 requires a
target holding cards, after which every remaining set must be declared on the log as it then stands.
The strategic content of the request survives; the endgame framing does not \cite{concession,rules}.
Measured incidence over 200 games: a team empties with sets unresolved in $41\%$ of games, $72\%$ of
those caused by an ask the other side could have declined to make --- but only $\mathbf{1.32}$ sets
remain unresolved on average when it happens, the clinch at five sets making deep endgames rare by
construction.

\begin{table}[t]
\centering
\caption{Stalling swept over trigger width and strength \cite{concession}. ``Decisions
changed'' counts trigger positions at which the mechanism moved the argmax.}
\label{tab:stall}
\footnotesize
\begin{tabular}{@{}rrrrrr@{}}
\toprule
trigger width & strength & trigger positions & decisions changed & win rate & paired diff \\
\midrule
$2$ & $1$   & $332$     & $16$              & $50.00\%$ & $+0.010 \pm 0.024$ \\
$6$ & $1$   & $1{,}952$ & $514$             & $49.67\%$ & $0.000 \pm 0.052$ \\
$3$ & $20$  & $552$     & $440$ ($80\%$)    & $42.00\%$ & $\mathbf{-0.42 \pm 0.12}$ \\
$6$ & $20$  & $1{,}471$ & $1{,}330$ ($90\%$) & $\mathbf{33.33\%}$ & $\mathbf{-1.14 \pm 0.19}$ \\
$6$ & $100$ & $1{,}421$ & $1{,}329$ ($94\%$) & $31.83\%$ & $-1.22 \pm 0.19$ \\
\bottomrule
\end{tabular}
\end{table}

Table~\ref{tab:stall} is inert when weak, severely harmful when strong, and monotone --- and not for
want of alternatives: at trigger positions there are on average $23.15$ legal asks across $2.71$
distinct targets, and only $5.9\%$ of positions offer a single target. The mechanism has room to
act, and acting is bad. Refusing a card you can certainly take refuses free material \emph{and} a
retained turn, and what it buys is time on a channel whose product is worth very little --- declares
are already about $100\%$ correct while the channel is open, and only about $1.32$ sets remain when
it closes. It is paying a certain card for insurance against a rare, small loss.
Section~\ref{sec:structural}'s caveat applies here with more force: a stall is a deliberately worse
ask, and the score cannot express one.

\section{The turn handoff: built, and the shipped rule already wins}
\label{sec:handoff}

The owner also asked that a seat which runs out of cards on its own turn choose the teammate to pass
to by reviewing the history --- which teammate could best use the turn --- rather than by hand size.
The threat layer of Section~\ref{sec:cost} is exactly that estimator, pointed at a teammate instead
of an opponent.

A win-rate A/B has no power at this decision, so the instrument is local: at every real pass
decision, fork the game once per candidate teammate, force the pass there, and count the cards that
teammate then actually takes before losing the turn --- ground truth per candidate
\cite{concession}. First, how often the decision exists at all, over 300 \usff{} games: $11.3\%$ of
games reach \code{awaitPass}; of those events $62.9\%$ present a real choice (more than one teammate
holding cards), $0.0\%$ have no teammate holding cards, and $11.4\%$ have the ask channel already
dead --- the owner's blind-declare case, which is 4 occurrences per 300 games. Mean unresolved sets
at the pass: $1.74$.

\begin{table}[t]
\centering
\caption{The fork study: 4{,}000 games, 223 real pass decisions, 446 candidate rows
\cite{concession}. ``Picks best'' is against the per-candidate ground truth the fork provides.}
\label{tab:handoff}
\footnotesize
\begin{tabular}{@{}lrrr@{}}
\toprule
estimator & corr.\ with cards taken & mean cards captured & picks best \\
\midrule
oracle & --- & $1.251$ & $100\%$ \\
\code{handSize} (shipped, \code{passTarget: 'most'}) & $\mathbf{-0.261}$ & $\mathbf{0.897}$ & $\mathbf{84.8\%}$ \\
\code{reach} (the threat estimator) & $\mathbf{+0.352}$ & $0.906$ & $85.2\%$ \\
hybrid (reach, hand size breaking ties) & --- & $0.892$ & $84.3\%$ \\
\bottomrule
\end{tabular}
\end{table}

Table~\ref{tab:handoff}'s correlations keep the signs of Section~\ref{sec:estimators} --- hand size
negative, reach positive --- and the \emph{decision} is still unaffected, because $\mathbf{73.1\%}$
of these decisions have no signal in them at all: every candidate takes the same number of cards, so
no estimator can separate them. Among the roughly 60 that do carry signal, hand size and reach
mostly agree. The shipped rule is already within $0.35$ cards of the oracle and picks correctly
$84.8\%$ of the time.

Capability~5 is therefore \textbf{built and refuted as a change}, not skipped. \textbf{The estimator
is right and the decision is empty}: $+0.352$ against $-0.261$ says the model is the better one, so
this is a statement about how rarely \usff{} asks the question, not a failure of the threat model.
\textbf{Where it would matter is frozen}: the states in which this decision fires in $41\%$ of games
belong to the 48-card rule set, which this project holds byte-identical \cite{styles}. And
\textbf{one premise of the request does not hold under \usff{}} --- ``which teammate could best
declare'' is not a reason to pass, because row 11 gives every seat the declare option without the
turn. The exception the owner identified is real and is the one case where it inverts: when the
opponents are all out, the declare option sits on the turn-holder alone --- the $11.4\%$ of pass
events above, or 4 occurrences per 300 games.

\section{Concealment: a counter-weapon, not an improvement}
\label{sec:conceal}

The owner asked for a fourth mechanism after the first three were measured. When an opponent asks
you for a card of set $H$, the natural reply is to ask into $H$ --- but that reply publishes your
own row-6 basis in $H$. The proposal is to hold the card quietly instead, so opponents infer you
have nothing there and skip you, until you can locate the rest of $H$ and declare it.

\subsection{The behaviour it targets, and why the mechanism is selectable at all}
\label{sec:arbitrate}

Before building a countermeasure, the behaviour was measured. After a seat is asked for a card of
$H$, misses, and takes the turn under row 10, $\mathbf{88.1\%}$ of its next asks go into $H$,
against an availability-matched control of $\mathbf{21.4\%}$ \cite{asking}. The pattern is real and
large. \textbf{But the engine is not reciprocating; it is collecting certainties.} $99.1\%$ of those
replies are the argmax on hit probability and $60$--$94\%$ are \emph{certain hits}: the opponent's
ask, combined with what the seat knew, located a card exactly, and the certainty bonus plus row 9 do
the rest. Concealment therefore proposes to decline a certain hit, the arithmetic that killed
stalling.

It nonetheless escapes Section~\ref{sec:structural}'s trap, because it is not a deliberate miss: it
moves the argmax \emph{between two hits} --- the concealed ask and some other set's ask --- and its
charge is paid at most once per set rather than at every decision, since a second ask into $H$
publishes nothing the first did not. So it is selectable at moderate weight, and measures non-zero
where signalling and stalling measured flat. It is charged as $a_c \cdot w_{\mathrm{hit}} \cdot c
\cdot \mathrm{mine}(H) / (1+E)$ --- deliberately the same conversion Equation~\eqref{eq:defuse}
uses, so that the two mechanisms arbitrate exactly rather than by accident. Setting the two against
each other, defusal wins iff $a_d \cdot p \cdot \mathrm{prey}(H) > a_c \cdot \mathrm{mine}(H)$, and
at both appetites $1$ that is simply
\begin{equation}
\label{eq:arbitrate}
p \;>\; \frac{\mathrm{mine}(H)}{\mathrm{prey}(H)}:
\end{equation}
take the card back when it is likely enough relative to how much of the set you would be
advertising. Measured over $13{,}622$ ask decisions the two mechanisms collide on $28.6\%$, and
defusal takes $79.8\%$ of the collisions. They compose; neither cannibalises the other.

\subsection{The measurement, including the part that decides it}
\label{sec:concealmeasure}

\begin{table}[t]
\centering
\caption{Concealment added to defusal, against opponents that defuse \cite{concession}. The
last row is the cell that decides what the mechanism is; Section~\ref{sec:floor} reclassifies it as
unresolved rather than zero.}
\label{tab:conceal}
\footnotesize
\begin{tabular}{@{}lrr@{}}
\toprule
arm & $N$ & paired set-difference \\
\midrule
byte-exact control & --- & $\mathbf{0.0000 \pm 0.0000}$ \\
held-out bank A & $800$ & $+0.8875 \pm 0.2443$ \\
held-out bank B & $800$ & $+0.8325 \pm 0.2378$ \\
headline & $600$ & $\mathbf{+0.9683 \pm 0.2744}$ \\
information-free null of the same shape and magnitude & --- & $\mathbf{-0.7200 \pm 0.2715}$ \\
against opponents that do NOT defuse & $600$ & $\mathbf{-0.1483 \pm 0.2571}$ \\
\bottomrule
\end{tabular}
\end{table}

Table~\ref{tab:conceal} carries a dose-response in how well the opponent can read the leak
concealment hides: $+0.97$ against a hard-skill opponent, $+0.31$ against medium, $+0.11$ against
one without constraints. \textbf{The entire effect is denial of defusal} --- not a general
improvement but a counter to one specific mechanism.

\textbf{A correction, printed where the claim was.} This section used to end with \emph{``removing
only the opponent's defusal drops the marginal value to $+0.157$''}. That is withdrawn. It
contradicts the table's own last row in sign, and it does not reproduce: measured directly,
conceal-only against plain is $-0.2883 \pm 0.2640$, and as a difference-of-differences on a common
bank $-0.3067$. Both readings come out negative, and where the $+0.157$ came from could not be
reconstructed, so it is retracted rather than reinterpreted \cite{concession}.

\textbf{And the row that decides the mechanism is below the instrument's resolution.} $-0.1483 \pm
0.2571$ at $N = 600$ sits under a minimum detectable effect of about $0.38$ sets at that cell's own
per-pair SD (Section~\ref{sec:floor}), so concealment against a non-defusing opponent is \textbf{not
resolved here}: anything from a real loss of about a third of a set to a small gain is consistent
with it. What \emph{is} resolved, at $4{,}000$ pairs per cell, is that conceal-only is the weakest
of the four concession arms outright (Section~\ref{sec:triangle}).

\subsection{Verdict, and the correction it forces on defusal}
\label{sec:concealverdict}

\textbf{Not shipped on.} \code{conceal} is landed as an optional style field defaulting to absent,
so every roster vector and every shipped difficulty tier is byte-identical, and the tiers carry
\code{defuse: 0} and so have no customer for a counter to defusal.

This section used to force a correction on Section~\ref{sec:defuse}, printed in the body as a
warning that \emph{``defusal's $+1.50$ is an over-estimate against a concealing opponent''} which
would take ``roughly one set per duplicate pair back'' --- on the grounds that neither the gauntlet
(Table~\ref{tab:gauntlet}) nor the cross-engine replication (Table~\ref{tab:xplay}) ever put
defusal in front of an opponent that conceals. \textbf{That gap has since been measured, and it
closes the opposite way} \cite{concession}. Table~\ref{tab:triangle}'s C vs D edge is precisely the
missing cell, at $4{,}000$ duplicate pairs per edge on two independent bank sets: a conceal-only
opponent does not take a set off a defusing one, it \emph{loses} to it, by $-0.6797 \pm 0.1030$ and
$-0.5102 \pm 0.1051$ --- about half a set per pair, in defusal's favour, on both banks.

The worry was the right shape aimed at the wrong target. The figure near $0.9$ belongs to a
different comparison entirely: concealment \emph{added to} defusal on the same side, the DC vs D
edge at $+0.9630$ and $+0.8207$ (Table~\ref{tab:triangle}). That prices what a defusing team gains
by concealing as well --- a fact about its own configuration --- and says nothing about what a
concealing opponent takes off it. Defusal's headline stands as measured.

\section{The triangle, refuted, and what replaced it}
\label{sec:triangle}

This project published a claim about its own tactic space: that defusal and concealment made it
\emph{intransitive}, and that the resulting cycle was the reason an arbitration layer over
per-decision tactics would pay. Measured properly, there is no cycle, and the retraction and its
replacement are reported here as a first-class result.

\begin{table}[t]
\centering
\caption{All six edges of the $2\times2$ through one code path, $4{,}000$ duplicate pairs
per edge, with only the arm definitions changing; then reproduced by an independent agent on six
entirely fresh banks \cite{concession}. Arms: \textbf{P} plain, \textbf{D} defuse-only, \textbf{C}
conceal-only, \textbf{DC} both.}
\label{tab:triangle}
\footnotesize
\begin{tabular}{@{}lrr@{}}
\toprule
edge & first run & independent replication \\
\midrule
D vs P   & $+1.6225 \pm 0.1016$ & $+1.6595 \pm 0.1022$ \\
\textbf{C vs D} & $\mathbf{-0.6797 \pm 0.1030}$ & $\mathbf{-0.5102 \pm 0.1051}$ \\
P vs C   & $+0.1457 \pm 0.1030$ & $+0.1742 \pm 0.1030$ \\
DC vs D  & $+0.9630 \pm 0.1069$ & $+0.8207 \pm 0.1070$ \\
\textbf{P vs DC} & $\mathbf{-1.2985 \pm 0.1036}$ & $\mathbf{-1.3155 \pm 0.1046}$ \\
DC vs C  & $+0.3690 \pm 0.0987$ & $+0.3202 \pm 0.0987$ \\
byte-exact control & $0.0000 \pm 0.0000$ & $0.0000 \pm 0.0000$ \\
\bottomrule
\end{tabular}
\end{table}

\textbf{There is no cycle under either reading of ``the concealing corner'', and the loop is broken
by a wrong-signed leg rather than a straddling one --- so no sample size rescues it.} Read the
corner as conceal-only and the second leg reverses: concealment does not beat defusal, it loses to
it by about half a set ($9.5$--$12.9$ SE from zero). Read it as defuse-plus-conceal --- which is
what Table~\ref{tab:conceal}'s $+0.97$ actually measured, concealment \emph{added to} defusal
against defusing opponents --- and the second leg holds but the third reverses instead: DC beats P
by $1.3$ sets ($12.5$--$24.6$ SE). The full $4\times4$ is a \textbf{strict transitive total order,
$DC > D > P > C$}, every edge the same sign on both bank sets. Concealment is an increment on
defusal, never a substitute for it, and \textbf{conceal-only is the worst of the four arms} ---
worse than plain.

Two further corrections came out of the replication rather than out of the run. The triangle table
used to quote defusal at $+1.92 \pm 0.31$ from a 400-pair common-baseline bank; at $4{,}000$ pairs
it is $+1.62$ / $+1.66$, ordinary sampling variation --- and the first explanation offered for that
gap, that the mirror harness lacks \code{pickAsk}'s certain-hit gate, was tested by running the
ungated harness and \emph{refuted} ($+1.6850 \pm 0.3358$, indistinguishable from the gated value).
And the claim that the published triangle mixed two definitions of the concealing corner is
\textbf{not} established: \code{conceal.ts}'s own header points the other way, and the two readings
cannot be told apart at this resolution.

\subsection{What survives, and why it is stronger than the cycle was}
\label{sec:onebit}

The old argument ran: the v1.0 adaptive layer degenerates because its counter table has a dominant
row \cite{v10paper}; a defuse/conceal pair has no dominant row; therefore reading the opponent
starts to pay. The premise was wrong. The conclusion is right anyway, for a plainer reason ---
\textbf{the column argmax is opponent-dependent.} Measured on common banks, paired on the deal twice
over, the best response against a \emph{plain} opponent is D by $+0.2870 \pm 0.1392$ over DC;
against a \emph{conceal-only} opponent it is again D, by $+0.1462 \pm 0.1437$ over DC; and against a
\emph{defuse-only} opponent it is DC, by $+0.9152 \pm 0.1353$ over D \cite{concession}.

The best response flips on one bit of the opponent's configuration --- \emph{does it defuse?} ---
and knowing that bit is worth up to $+0.29$ sets per duplicate pair against a non-defusing opponent,
which is not a hypothetical opponent: every shipped tier carries \code{defuse: 0}. It is
Section~\ref{sec:concealmeasure}'s own thesis restated as a policy --- conceal only against an
opponent that defuses --- and it makes the follow-on work concrete and \emph{smaller} than a
portfolio. The v1.0 degeneracy verdict does need revisiting, since degeneracy there is defined as
``the read never changes the argmax'' and here the read changes the argmax; but the table it needs
is one binary coordinate, not a tenth style \cite{v10paper,adaptive}. The arbitration layer is still
the larger question and still unbuilt.

\section{The detection floor, and the two nulls it reclassifies}
\label{sec:floor}

Every headline in this paper comes from one instrument: $K$ duplicate pairs of arm $X$ against arm
$Y$ under \usff{}, reported as the mean paired set-difference $\pm 1.96$ SE. Until it was
characterised, every published null was ambiguous between ``no effect'' and ``no resolution''. That
characterisation is a sibling result in this batch \cite{floorpaper}; what matters here is what it
does to this paper's own numbers. The per-pair SD is $3.15$--$3.44$ sets for arms that differ at
every decision (replication: $3.18$--$3.44$), giving an $80\%$-power minimum detectable effect of
$0.76$ at $N = 150$, $0.54$ at $300$, $0.47$ at $400$, $0.38$ at $600$, $0.33$ at $800$, $0.21$ at
$2{,}000$ and $0.14$ at $4{,}300$ pairs.

The power model is the solid part: across twelve (planted-effect, $N$) cells the observed detection
rate sat within about one binomial SE of normal-theory prediction in both runs. Two parts did not
replicate, and both cut against the first run's conclusions --- the empirical floor at $N = 400$
($84.7\%$ detection of a planted $0.393$, then $75.0\%$; pooled $81.7\%$, so the floor is somewhere
in $0.36$--$0.50$) and interval coverage at $N = 400$ ($6.1\%$ false positives, then $10.8\%$;
pooled $8.0\%$ $[5.2\%, 11.7\%]$, excluding the nominal $5\%$). At $N = 800$ both runs are nominal
\cite{floorpaper}.

\begin{table}[t]
\centering
\caption{Which numbers in this paper sit below their own cell's floor
\cite{concession,floorpaper}. The MDE depends on \emph{that cell's own} SD; this is a record of
specific verdicts, not a threshold to apply by eye.}
\label{tab:floortable}
\footnotesize
\begin{tabular}{@{}lrrrl@{}}
\toprule
result & reported & $N$ & its floor & verdict \\
\midrule
Table~\ref{tab:conceal}, conceal vs non-defusers & $-0.1483 \pm 0.2571$ & $600$ & ${\approx}\,0.38$ & \textbf{below floor --- not resolved} \\
The contained-pass aim, ``inert'' & $-0.045 \pm 0.130$ & --- & --- & \textbf{below floor --- a non-measurement} \\
\S\ref{sec:standing}, all four headline cells & $42.7$--$48.7\%$, no CIs & $150$ & $1.05$--$1.13$ & \textbf{all four below floor} \\
Table~\ref{tab:defuse}, defusal shipped & $+1.43$ / $+1.58 \pm 0.32$ & $400$ & ${\approx}\,0.47$ & resolved, $3\times$ the floor \\
Table~\ref{tab:triangle}, the six edges & $\pm 0.10$ half-widths & $4000$ & ${\approx}\,0.15$ & resolved \\
\bottomrule
\end{tabular}
\end{table}

Table~\ref{tab:floortable} is the discipline applied to this project's own work. \textbf{Two of this
paper's nulls are reclassified as non-measurements}: the contained-pass aim called ``inert'' in
Section~\ref{sec:defects}, and the non-defuser cell of Table~\ref{tab:conceal}. Neither is a
statement about the game; both are statements about this instrument's resolution. A third
reclassification lands on the cross-engine standing figures: at 150 games per cell the per-opponent
ranking of Section~\ref{sec:standing} is not resolved, which is why Section~\ref{sec:crossplay}
rests on Table~\ref{tab:xplay}'s 300 properly paired duplicate pairs and not on it.

One trap is worth naming, because the first floor run fell into it: the MDE depends on that cell's
own SD, and the SD is not constant across the repository --- the turtle cell of
Table~\ref{tab:defuse} measures $3.645$, and the cross-engine harness's own published half-widths
imply it runs at $4.57$--$4.95$, some $38$--$49\%$ higher \cite{floorpaper}. So one SD substituted
across all conditions understates the floor exactly where the cells are smallest. Everything this
paper ships a mechanism on --- Tables~\ref{tab:defuse}, \ref{tab:gauntlet} and~\ref{tab:triangle}
--- clears its own floor by $3\times$ or more.

\section{Two defects in the shipped engine, found on the way}
\label{sec:defects}

\textbf{1. \code{missTarget: 'fewest'} is backwards.} Its trace prose tells the user that the
smallest hand makes the surrendered turn ``worth least''; the measurement says the opposite in every
phase bucket (Section~\ref{sec:estimators}), and the fork study of Section~\ref{sec:handoff}
reproduces the sign a third time at $-0.261$. The same proxy appears in \code{valueContainedPass}'s
$n_t$ term and in the contained-book policy's aiming rule \cite{containedpaper}.

\textbf{The fix was built where it is soundest, and measured inert.} A contained-book turn-pass is a
\emph{guaranteed miss} by construction, so Section~\ref{sec:confound}'s danger/opportunity confound
cannot arise there: it is the one place in the engine where ``concede to the least dangerous seat''
is exactly right. Replacing both the aim and the cost model with the threat estimate measured
$-0.045 \pm 0.130$ over 400 duplicate pairs, and the diagnostic says why: over $25{,}654$ ask
decisions, more than one opponent still holds cards $99.1\%$ of the time, the contained pass fires
at all on $\mathbf{1.06\%}$, and it fires with the two aims disagreeing on $\mathbf{0.11\%}$.
Published reach is too sparse there: most opponents have shown no basis at all, so every candidate
prices the same and the aiming gain collapses.

\textbf{The change was reverted rather than shipped off by default}, on the argument this project
applies to its own knob inventory --- a knob that is listed, swept and reported while being inert
contaminates the measured vector, and one of those is already one too many. Section~\ref{sec:floor}
then reclassifies the $-0.045$ itself: it is below the floor, so ``inert'' overstates what was
measured, and the revert stands on the sparsity diagnostic and the contamination argument rather
than on the null.

\textbf{2. \code{style.signalling} is unreachable under \usff{}.} The signalling ask is gated inside
the planner path, which a \usff{} game enters only for \code{phase === 'endgame'} --- a phase that
rule set never produces. So the knob is listed, swept and reported, and buys a \usff{} seat nothing
at all except through the contained-pass information price, where it zeroes one term. \textbf{A dead
knob in a measured vector is a contaminated comparison}, the same argument that decided the revert
above. This project has published that failure mode once before, for a different axis
\cite{inertpaper}; this is the second instance in the same roster.

\section{Capability 4: switching actions instead of playstyles}
\label{sec:capability4}

The owner asked for an adaptive model that, ``instead of switching between playstyles, just switches
between actions and strategies, to allow for more flexible playing''. \textbf{That is the right
architecture, and it is now measured rather than argued.} The two layers price against each other in
the same unit: the v1.0 adaptive layer, which switches the whole style once per phase, is worth
$+0.13 \pm 0.06$ sets per duplicate pair over a bare Balanced team \cite{v10paper,adaptive,bounded};
the style knob itself, chosen well, is worth nothing, because it is a constant (a $\pm 0.27$
spread); and the per-decision terms (\code{containedPass}, \code{defuse}), which switch the
\emph{action} at every decision from the position, are worth $+1.43$ and $+1.58 \pm 0.32$
(Table~\ref{tab:defuse}).

\textbf{The per-decision layer is worth roughly eleven to twelve times the style-switching layer},
and the style-switching layer is worth less than simply naming a better style would have been ---
the owner's point, quantified. One arithmetic correction belongs here rather than in a footnote: the
project's own summary table records this ratio as ``about $15\times$'', and against the numbers this
paper reports it is not, since $1.43/0.13 = 11.0$ and $1.58/0.13 = 12.2$. No source records what
the larger figure was computed from. The plausible reconstruction --- offered as one, not as a
provenance --- is that it is the ratio against the defusal estimate this paper has withdrawn,
Section~\ref{sec:triangle}'s retired $+1.92$, since $1.92/0.13 = 14.8$. What is not in doubt is the
ratio of the numbers actually carried, which is eleven to twelve.

There is a structural reason the style layer cannot do better, and it is a theorem rather than a
measurement: best-response selection over the committed counter table provably degenerates to
always-Punter, because one row weakly dominates every column and the expected payoff is linear in
the opponent posterior \cite{v10paper,adaptive}. A layer that reads the \emph{position} --- has this
opponent published a basis in a set my team can locate cards of? is this set contained? is the ask
channel about to close? --- is choosing from the legal move list instead, every turn, and every
mechanism in this paper is one predicate over that list. So what is shipped is a \emph{step} of
capability 4, and the arbitration layer over it is worth building only now that there is more than
one tactic to arbitrate: of the four originally proposed three measured harmful or inert, and
Section~\ref{sec:onebit} prices the arbitration between the two survivors at one bit.

\section{Threats to validity}
\label{sec:threats}

\begin{enumerate}
\item \textbf{One fitted constant, fitted here.} The cards-per-prey-card slope of
  Equation~\eqref{eq:threat} is measured on this engine, against this roster, under \usff{}.
  Table~\ref{tab:xplay} is evidence that the \emph{mechanism} transfers, not that the constant is
  optimal elsewhere; no re-fit was attempted inside the host engine, because that would have burned
  it as a holdout \cite{crossplay}.
\item \textbf{A stated $K$ is missing.} Several variants of the defusal term were scored before the
  winner was picked --- per-set against per-seat, five appetites, two gate designs --- and the
  selection haircut is unpriced. The held-out banks and the cross-engine replication stand in for
  it; neither substitutes for the arithmetic.
\item \textbf{The threat model is an under-approximation by construction.} It reads no hand but the
  viewer's own, so it is silent about every basis a seat has not published and is worst early. That
  property is what makes avoidance unsound (Section~\ref{sec:confound}); it also caps how much of
  the true threat any of these mechanisms can act on.
\item \textbf{The opponents are all policies of two engines}, and both populations are closed. The
  gauntlet's opponents are FishAI policies; the cross-engine opponents are four bots roughly four
  releases behind their own project's frontier, played without the declare channel that carries
  $45.4\%$ of FishAI's declares (Section~\ref{sec:bridge}). Neither population contained a
  \emph{concealing} opponent --- but that particular gap is closed rather than open:
  Table~\ref{tab:triangle}'s C vs D edge measures defusal against a conceal-only opponent at
  $4{,}000$ pairs on two bank sets and puts defusal ahead by about half a set
  (Section~\ref{sec:concealverdict}).
\item \textbf{A channel the host engine has and FishAI does not was switched off.} Their engine
  supports \code{psychologicalTells}. FishAI has no notion of it, so it was left \textbf{off} for
  every cell in Section~\ref{sec:crossplay}, on the grounds that leaving it on would have measured
  something other than card play \cite{crossplay}. That is the right call for isolating the
  mechanism, and it is also a bound on the standing figures: their bots are not being played with
  everything the host gives them.
\item \textbf{Resolution, and the harnesses.} Section~\ref{sec:floor} is the honest boundary of
  every null above, and a null not named in Table~\ref{tab:floortable} has not been individually
  checked against its own cell's SD. The measurements come from scratch harnesses under
  \code{scripts/probe-*.mjs}: not typechecked, no committed artifact, not digest-bound, unlike the
  lab's tournament artifacts \cite{botlab,repo}. They exist so every table above can be regenerated,
  and that is the strongest claim available for them.
\end{enumerate}

\section{What is not settled}
\label{sec:open}

The residue is kept in the body, because a paper that shipped a mechanism and hid its open items
would be misreporting the work.
\begin{itemize}
\item \textbf{The style matrix is invalidated.} The roster now carries \code{defuse: 1}, so the
  committed style matrix no longer describes the shipped roster and must be re-run and re-pinned
  \cite{styles}.
\item \textbf{A pre-existing defect in the committed search numbers}, found while adding the new
  knobs to the ladder: the candidate budget was the literal $60$ against a $75$-rung ladder, so the
  search evaluated a \emph{prefix} and abandoned the tail --- a prune selected by the order the
  fields happen to be typed in. Replaying the old ladder along each style's committed accepted moves
  reproduces the artifact's candidate counts exactly and shows what the prefix cost: \code{banker}
  never evaluated \code{minHandSize=3}, and \code{turtle} never evaluated \code{minHandSize=1} or
  \code{minHandSize=3} --- three rungs across two of nine searches. The budget is now the ladder's
  own length. The exploit cache does not see any of this, because it keys on an engine hash that
  walks the engine tree only and neither the ladder nor the budget lives there, so a cached re-run
  re-emits the pre-fix numbers.
\item \textbf{The concealment half of the certain-hit gate is unpinned.} \code{pickAsk} documents
  the gate as zeroing \emph{both} terms; the gate test pins only defusal, because in its fixture the
  concealing seat has already published a basis in both sets it can ask into and the charge is
  exactly zero on both asks --- asserted explicitly, so the test cannot be misread as coverage it
  does not have. A mutation leaving that charge live under a certain hit is invisible to the suite,
  and the gate's own firing rate is unpriced by any ladder (Section~\ref{sec:gate}).
\item \textbf{Sign-flip coverage was over-claimed and then corrected.} The claim was that both
  concession terms pass their suites with the sign reversed. Measured by reversing each of the three
  sign errors available: negating the defusal bonus fails 2 of 25 existing assertions, negating the
  concealment charge fails 3 of 25, and reversing the \emph{application} in \code{pickAsk} ---
  writing $\text{charge} - \text{bonus}$ for $\text{bonus} - \text{charge}$ --- \textbf{passes both
  suites entirely}. One third of the claim was true and it was the important third: every assertion
  called the exported functions directly, and nothing pinned what \code{pickAsk} did with the
  numbers they returned. Two decision-level tests now close it.
\item \textbf{Four seams are held by hand.} No test pins the concession layer's turn-yield equal to
  the contained-book policy's $E$, though both file headers describe them as deliberately
  duplicated; the tabled-licence export has no caller; the dynamic test of the public-view boundary
  runs the three shipped tiers, which carry \code{defuse: 0}, so no code in the threat, defusal or
  concealment files executes under it; and the bounded arm's licence scan retires licences against
  budgeted knowledge while the ask history it scans is not itself budgeted, a fix that would change
  every committed v1.5 number \cite{v15paper,bounded}.
\item \textbf{A true match against SESTINA v1.0 is possible and is blocked only by a missing
  compiler} (Section~\ref{sec:sestina}). Their repository documents a JSON-line cross-play protocol;
  installing a C++ toolchain would make it a day's work \cite{crossplay,fishlabs}.
\end{itemize}

\section{Conclusion}
\label{sec:conclusion}

The request that started this work was for a rule that refuses to ask dangerous opponents. The rule
was built exactly as asked, measured against a byte-exact control on duplicate deals, and it loses
--- up to $11.7$ points of win rate, monotonically worse as it is driven harder, while an
information-free perturbation of the same magnitude is free. The threat model underneath the request
is nonetheless correct, and better than what the engine already used for the same job: the shipped
hand-size proxy correlates $-0.147$ with what a conceded seat actually takes, the licence estimator
$+0.421$.

What makes the prescription wrong is one structural fact. Row 6 permits an ask only into a set the
asker holds a card of, so the seat that is dangerous to you is the seat holding cards of the sets
you are invested in --- which is the seat most likely to hold the card you want. Threat and
opportunity are the same measurement, $\mathrm{corr} = +0.249$, and avoidance trades away your own
best asks to buy safety that a public-log threat model can never certify anyway, since an absence of
published reach is not evidence of absent reach. Invert it and the same model prescribes the
opposite move: ask \emph{into} the set the dangerous seat has published a basis in, take the card
its reach rests on, and keep the turn under row 9. That is worth $+1.43$ and $+1.58$ sets per
duplicate pair on held-out banks, is positive at all nine roster styles, survives opponents that
defuse back at $+1.08$ to $+1.96$, and --- the result the project would keep if it could keep only
one --- transfers at the same magnitude, $+1.55$ to $+1.75$, into a third party's engine against a
third party's bots with one knob moved.

The rest of the ledger is negative and is reported at the same weight: signalling flat and
unreplicated, stalling harmful once strong enough to act, the turn-handoff request beaten by the
rule already shipped, concealment a counter-weapon rather than an improvement. Three items in it are
this project correcting itself. The intransitive triangle it published about its own tactic space
does not exist, and the opponent-dependent argmax worth one bit that replaced it is the better
result. A fix for the engine's wrong-signed concession proxy was built where it is provably sound,
measured below the instrument's floor, and \emph{reverted rather than shipped off by default},
because a listed-but-inert knob contaminates every vector it appears in. And the certain-hit gate
exists because an adversarial review found the new credit could displace a certain hit, which it
then did nine times in a sweep against zero for the control.

Two of the nulls above are not zeros. The detection floor \cite{floorpaper} reclassifies
concealment's non-defuser cell and the contained-pass aim as \emph{non-measurements} --- this
project's discipline applied to its own published numbers rather than to someone else's. Where a
result is below the instrument's resolution, the honest report is the resolution.

\medskip
\noindent\textbf{Reproducibility.} The engine is deterministic and one seed reproduces a
byte-identical game \cite{repo}; every table above is regenerated by a named probe under
\code{scripts/probe-*.mjs} \cite{concession}. The cross-engine harness lives outside the repository,
beside the clone it drives: that clone carries no licence file, nothing from it is copied into this
project, and their findings are cited as theirs \cite{crossplay,fishlabs}.

\begin{thebibliography}{99}
\footnotesize

\bibitem{concession}
FishAI project. CONCESSION.md --- FishAI v2.0: the three-sided ask. Repository document, 2026. Every
measurement in this paper except the cross-engine ones, with its probe named: \S1--\S3 the conceded
turn, off-limits and defusal, \S4--\S5a signalling, stalling and concealment, \S6--\S8 the engine
defects, the turn handoff and the triangle, \S8a the detection floor, \S9 the open residue.

\bibitem{crossplay}
FishAI project. CROSSPLAY.md --- FishAI against the fishlabs bots, in the fishlabs engine.
Repository document, 2026. \S0 what ran and what did not, \S1 the bridge and the declare-window
handicap, \S2 standing, \S3 the cross-engine defusal replication, \S4 their reported SESTINA
results and their own qualifications on them, \S5 \code{lockout}, \S6 common failure points and
declare calibration, \S7 what it does not establish.

\bibitem{asking}
FishAI project. ASKING.md --- the ask channel and what it publishes. Repository document, 2026.
\S4.3 the reply-into-the-asked-set rate and its availability-matched control.

\bibitem{containment}
FishAI project. CONTAINMENT.md --- the contained-book result. Repository document, 2026. Claims
C1--C6 with evidence tiers, and the independent derivation of the value of a turn.

\bibitem{adaptive}
FishAI project. ADAPTIVE.md --- the v1.0 adaptive layer and its degeneracy verdict. Repository
document, 2026. The dominant-row argument that collapses best-response style selection to a
constant.

\bibitem{bounded}
FishAI project. BOUNDED.md --- the bit-budget memory ladder. Repository document, 2026. \S5a the
v1.0-over-Balanced measurement cited in Section~\ref{sec:capability4}.

\bibitem{rules}
FishAI project. RULES\_US54.md --- the ``US student'' 54-card variant. Repository rules
specification with derived consequences and test vectors, 2026. Decision-table rows cited by number
throughout.

\bibitem{styles}
FishAI project. STYLES.md --- the play-style roster under \usff{}. Repository document, 2026. The
nine-style roster, the rule that a style may not tune its own trigger, and the byte-identical hold
on the 48-card rule set.

\bibitem{botlab}
FishAI project. BOT\_LAB.md --- play-style spectrum: research, experimental design, and metrics.
Repository document, 2026. \S5.1 duplicate-deal pairing, \S10 threats to validity.

\bibitem{repo}
FishAI project. FishAI repository: engine, bots, simulation laboratory, and results site. Concession
layer \code{lib/engine/bots/threat.ts}, \code{defuse.ts} and \code{conceal.ts}, the ask ranker
\code{decide.ts} (\code{pickAsk}, the certain-hit gate), and
\code{tests/bots/concession-direction.test.ts}, 2026.

\bibitem{v05paper}
A.~Hsieh. FishAI v0.5: measuring play style without skill in a 54-card Literature variant. FishAI
project; companion paper, 2026.

\bibitem{v10paper}
A.~Hsieh. FishAI v1.0: when best-response adaptation is worth less than nothing. FishAI project;
companion paper, 2026.

\bibitem{v15paper}
A.~Hsieh. FishAI v1.5: a bit-budget memory ladder as an honest difficulty axis. FishAI project;
companion paper, 2026.

\bibitem{obspaper}
A.~Hsieh. What a public log reveals: observing play style in Literature (Canadian Fish). FishAI
project; companion paper, 2026.

\bibitem{inertpaper}
A.~Hsieh. The inert axis: when a style parameter is wired, swept, and never reached. FishAI project;
companion paper, 2026. The first instance in this roster of the failure mode
Section~\ref{sec:defects} reports a second one of.

\bibitem{containedpaper}
A.~Hsieh. The contained book: an absorbing resource measured to be worth nothing. FishAI project;
companion paper, 2026. The value of a turn, the trigger inequality, and the aiming rule
Section~\ref{sec:defects} attempted to correct.

\bibitem{floorpaper}
A.~Hsieh. The detection floor: what a duplicate-deal instrument can and cannot resolve. FishAI
project; companion paper, published alongside this one, 2026. The per-pair SD, the power model, the
two parts that did not replicate, and the reclassification of two of this paper's nulls.

\bibitem{fishlabs}
\code{dylann4500}. FishLab --- a Canadian Fish engine, bot family and evaluation lab. GitHub
repository, \url{https://github.com/dylann4500/FishLab}. The C++ SESTINA v1.0 agent and its
reported margins, the TypeScript lab engine and its eight bots, the \code{lockout} danger-penalty
strategy, and the JSON-line cross-play protocol. The repository carries no licence file; nothing
from it is reproduced here.

\bibitem{pagat}
J.~McLeod (ed.). Literature (Canadian Fish). Pagat card-game rules archive,
\url{https://www.pagat.com/quartet/literature.html}. Folklore source for the common core of both
rule sets this project supports.

\end{thebibliography}

\end{document}
